\chapter*{Résumé}
\addcontentsline{toc}{chapter}{Résumé}

\paragraph{Pourquoi une version 0.7.} L'audit de la version 0.6 a révélé un défaut structurel : un ralentissement de la
productivité ou une baisse de la population active y réduisaient le chômage \emph{et} la dette publique. La cause est
l'absence de bloc d'offre : sans PIB potentiel, le modèle ne distinguait pas une baisse de la demande d'une baisse de
la capacité de produire. La v0.7 ajoute un PIB potentiel, une fermeture sur l'écart de production avec anticipation, et
dix-huit corrections ou mécanismes comportementaux (offre de travail calculée avec OpenFisca, transfert de bénéfices,
répercussion asymétrique de la TVA, épargne de précaution, rigidité des salaires, stabilisateur chômage, délocalisation
énergétique, R\&D, formation). Chaque paramètre est sourcé, avec une fourchette basse et haute.

\paragraph{Validation.} Les chocs d'offre ont désormais le bon signe : $-0{,}3$ point de productivité par an réduit le
PIB de 3,2 \% et augmente la dette de 1,8 point en 2037 (v0.6 : $-0{,}9$ \% et $-4{,}5$ points). Une baisse de 20\md\ de
cotisations crée 120 000 à 160 000 emplois, en accord avec les évaluations du CICE. Douze tests automatisés de
cohérence comportementale passent. Le choc de productivité n'est pas validé sur données historiques : c'est une
cohérence théorique.

\paragraph{Ce qui change par rapport à la v0.6.} Les mesures fiscales et budgétaires gardent des effets proches. Les
réformes d'offre changent nettement : relever l'âge légal à 65 ans réduit désormais la dette (1,6 point en 2037) au lieu
de l'accroître ; la prime d'activité augmente la population active, le RSA la réduit légèrement ; une hausse d'IS
rapporte 40 \% de moins.

\paragraph{Optimisation.} À partir de 35 leviers, cinq stratégies ont été construites, chacune par un objectif et des
contraintes explicites, puis vérifiées par le modèle complet, à 5, 10, 15 et 20 ans, sous dix environnements, trois
niveaux de comportement, deux intensités d'anticipation et une mise en œuvre partielle. Principaux enseignements :
\begin{itemize}
\item La stratégie de \textbf{redressement accéléré} (S4) réduit la dette de 19 points en 2037 et de 57 points en
2047, de façon robuste, au prix de 0,8 point de PIB, de 0,3 point de chômage en 2037 et d'une perte de niveau de vie
de 0,5 à 1,5 \% pour tous les déciles : acceptabilité faible.
\item La stratégie de \textbf{protection} (S1) augmente le niveau de vie de tous les déciles (de 7,8 \% pour le
premier à 2,2 \% pour le dernier), réduit la pauvreté de 3,5 points et le chômage de 1,4 point, sans accroître la
dette en 2037 ; elle repose sur une hausse de l'IS et des droits de succession dont le coût pour les ménages n'est pas
mesuré.
\item Les stratégies \textbf{emploi} (S2) et \textbf{productivité} (S3) ont les plus forts effets sur l'activité
(+2,4 et +4,4 \% de PIB en 2037), mais accroissent la dette pendant la transition et reposent sur des hypothèses
fragiles ; pour S3, les mesures d'offre ne s'additionnent pas (dette +6 points en 2037 au lieu de l'équilibre).
\item L'\textbf{optimisation sans pondération} (S5) est modeste mais robuste (+1,1 \% de PIB, $-1{,}7$ point de dette en
2037, $-10$ en 2047), avec des perdants limités.
\item Leviers retenus quelles que soient les préférences : prime d'activité, RSA, droits de succession, baisse de la
CSG sur les revenus d'activité, âge légal, investissement public.
\end{itemize}

\paragraph{Limites.} Effets distributifs statiques, entreprises non différenciées par taille, taux souverains exogènes,
anticipation non estimée, efficacité de la dépense publique non représentée. Elles sont détaillées au
chapitre~\ref{ch:limites} avec les travaux prioritaires.

\begin{hypothese}[Comment lire ce rapport]
Les résultats sont des \emph{écarts à une trajectoire de référence conventionnelle} (politique inchangée), et non des
prévisions. Chaque chiffre provient d'une simulation exécutée ; les fichiers sont fournis. Les niveaux de preuve des
paramètres sont indiqués : E (estimé sur données françaises), L (littérature), H (hypothèse), X (exploratoire).
\end{hypothese}
